\def\probtitle{같게 만들기}
\def\probno{H} % 문제 번호

\begin{problem}{\probtitle}{표준 입력(stdin)}{표준 출력(stdout)}{1 초}{1024 megabytes}

길이가 \(N\)인 정수 수열 \(A\)와 \(Q\)개의 쿼리가 주어진다. 각 쿼리는 다음과 같은 형태이다.

\begin{itemize}
    \item \(i~x\): \(A_i\)를 \(x\)로 변경한다.
\end{itemize}


초기 수열에 대해, 그리고 각 쿼리를 처리한 직후 수열에 대해 각 수열을 \textbf{같게 만들 수} 있는지 판별하라.

어떤 수열 $B$를 \textbf{같게 만들 수} 있다는 것은, 다음 연산을 0회 이상 사용해서 수열의 모든 원소를 같은 값으로 만들 수 있다는 의미이다.

\begin{itemize}[topsep=0pt,noitemsep]
    \item 두 정수 $i, j$를 고른다. $i$와 $j$가 같을 수 있다. ($1 \le i, j < N$)
    \item 그 후 수열을 다음과 같이 변경한다.
    \begin{enumerate}[topsep=0pt,noitemsep] 
        \item $B_i$를 $1$ 증가시킨다.
        \item $B_{i+1}$을 $1$ 감소시킨다.
        \item $B_j$를 $1$ 감소시킨다.
        \item $B_{j+1}$을 $1$ 증가시킨다.
    \end{enumerate}
\end{itemize}


% 길이가 $N$인 정수 수열 $A$가 주어진다. 당신은 다음 연산을 원하는 만큼 수행할 수 있다.

% \begin{itemize}[topsep=0pt,noitemsep]
%     \item 두 정수 $i, j$를 고른다. ($1 \le i, j < N$)
%     \item 그 후 수열을 다음과 같이 변경한다.
%     \begin{itemize}
%         \item $A_i$를 $1$ 증가시킨다.
%         \item $A_{i+1}$을 $1$ 감소시킨다.
%         \item $A_j$를 $1$ 감소시킨다.
%         \item $A_{j+1}$을 $1$ 증가시킨다.
%     \end{itemize}
% \end{itemize}

% 이후 아래의 쿼리가 $Q$개 주어진다.
% \begin{itemize}
%     \item $i~x$: $A_i$를 $x$로 변경한다.  
% \end{itemize}

% 초기 수열에 대해, 그리고 각 쿼리를 처리한 직후의 수열에 대해, 위 연산을 적절히 반복하여 수열의 모든 원소를 같은 값으로 만들 수 있는지 판별하여라.

\InputFile

첫째 줄에 수열 $A$의 길이를 나타내는 정수 $N$이 주어진다.

둘째 줄에 $N$개의 정수 $A_1, A_2, \cdots, A_N$이 공백으로 구분되어 주어진다.

셋째 줄에 정수 $Q$가 주어진다.

이후 $Q$개의 줄에 걸쳐 두 정수 $i, x$가 공백으로 구분되어 주어진다. 이는 $A_i$를 $x$로 변경하는 것을 의미한다.

\OutputFile

총 $Q+1$개의 줄을 출력한다.

첫째 줄에는 초기 수열을 \textbf{같게 만들 수} 있으면 \texttt{Yes}를, 그렇지 않으면 \texttt{No}를 출력한다.

그 다음 $Q$개의 줄에는 각 쿼리를 처리한 직후의 수열에 대해, \textbf{같게 만들 수} 있으면 \texttt{Yes}를, 그렇지 않으면 \texttt{No}를 출력한다.

\Constraints

\begin{itemize}[topsep=0pt,noitemsep]
    \item 주어지는 모든 수는 정수이다.
    \item $2 \le N \le 200\,000$
    \item $1 \le Q \le 200\,000$
    \item $1 \le j \le N$인 모든 $j$ 에 대해 $1 \le A_j \le 10^8$
    \item 모든 쿼리에 대해 $1 \le i \le N$,~ $1 \le x \le 10^8$
\end{itemize}

\Example

\begin{example}
    \exmpfile{./example/01.in.txt}{./example/01.out.txt}%
    \exmpfile{./example/02.in.txt}{./example/02.out.txt}%
\end{example}

\Notes
초기 수열이 $A = [3, 1, 1, 3]$일 때, $i = 3$, $j = 1$을 선택하여 연산을 한 번 수행하면
수열은 $A' = [2, 2, 2, 2]$가 된다. 따라서 초기 수열을 \textbf{같게 만들 수} 있다.

1번 쿼리로 $A_2 = 3$을 적용하면 수열은 $A = [3, 3, 1, 3]$으로 변경된다.
이 경우에는 수열을 \textbf{같게 만들 수} 없다.

2번 쿼리로 $A_3 = 3$을 적용하면 수열은 $A = [3, 3, 3, 3]$으로 변경된다.
이 경우에는 수열을 \textbf{같게 만들 수} 있다.
\end{problem}